\documentclass[10pt]{article}

\usepackage[utf8]{inputenc}

\usepackage[T1]{fontenc}

\usepackage{amsmath}

\usepackage{amsfonts}

\usepackage{amssymb}

\usepackage[version=4]{mhchem}

\usepackage{stmaryrd}



\title{$\mathbf{I}$ }



\author{}

\date{}





\begin{document}

\maketitle

ЯДЕРНАЯ $\boldsymbol{C}^{*}$-АЛГЕБРА - $C^{*}$-алгебра $A$, обладающая следующим свойством: для любой $C^{*}$-алгебры $B$ в алгебраическом тензорном произведении $A \otimes B$ этих алгебр существует единственная норма, пополнение в к-рой превращает $A \otimes B$ в $C^{*}$-алгебру. Таким образом, Я. $C^{*}$-а. по отнотениго к тензорным произведениям ведут себя подобно ядерным пространствам (хотя бесконечномерные Я. $C^{*}$-а. не являются ядерными пространствами). Класс Я. $C^{*}$-а. включает в себя все $C^{*}-$ алгебры типа 1. Этот класс замкнут относительно индуктивного предела. Если $I$ - замкнутый двусторонний идеал в $C$ *-алгебре $A$, то $A$ ядерна тогда и только тогда, когда ядерны $I$ и $A / I$. Подалгебра в Я. $C^{*}$-а. не обязана быть Я. $C^{*}$-а. Тензорное произведение $C^{*}$-алгебр $A$ и $B$ ядерно тогда и только тогда, когда $A$ и $B$ (обе) являются Я. $C^{*}$-а. Если $G$ - аменабельная локально бикомпактная группа, то обертывающая $C^{*}$-алгебра групповой алгебры $L^{1}(G)$ ядерна (обратное неверно). Каждое факторпредставление Я. $C^{*}-$. гиперфинитно, т. е. порождаемая этим представлением Неймана алгебра может быть получена из возрастающей последовательности конечномерных факторов (матричных алгебр). Любое факторсостояние на ядерной $C^{*}$-подалгебре в $C^{*}$-алгебре продолжается до факторсостояния на всей алгебре.



Пусть $L(H)$ есть $C^{*}$-алгебра всех ограниченных линейных операторов в гильбертовом пространстве $H, A$ есть $C^{*}$-алгебра операторов в $H$. Если $A-$ Я. $C^{*}$-а., то ее слабое замыкание $\bar{A}$ является и н ь е кт и в о й а лге б рой Н е й м а на, т. е. существует проекция $L(H) \rightarrow \bar{A}$ с единичной нормой; в этом случае коммутант $A^{\prime}$ алгебры $A$ также инъективен. Произвольная $C^{*}$-алгебра $A$ ядерна тогда и только тогда, когда ее обертывающая алгебра Неймана инъективна. $C^{*}$-алгебра $A$ ядерна тогда и только тогда, когда она обладает свойством в по лн е по ло жи т е л ь н о ї а I I р о к с и м а ц и и, т. е. тождественный оператор в $A$ аппроксимируется в сильной операторной топологии линейными операторами конечного ранга с нормой, не превосходящей 1 , и с нек-рым дополнительным свойством «вполне положительности» [1].



Јюбая Я. $C^{*}-\mathbf{a}$. обладает свойством аппроксимации и ограниченной аппроксимации (см. Ядерный оператор). Однако существует неядерная $C^{*}$-алгебра со свойством ограниченной аппроксимации. $C^{*}$-алгебра $L(H)$ всех ограниченных операторов в бесконечномерном гильбертовом пространстве $H$ не обладает свойством вполне положительной аппроксимации и даже свойством аппроксимации, так что $L(H)$ не является Я. $C^{*}$-а



Jum.: [1] L a n c e F. Ch., «Proc. Symp. Pure Math.», 1982, v. 38, pt 1, p. 379-99; [2] Б р а т т е л и У., Р о б и н с о н Д., Операторные алгебры и квантовая статистическая механика,

пер. с англ., М., 1982.

М. Литвинов, Я̆ЕРНАЯ 'ИИЛИНЕЙААЯ ФОРМА - билинейная форма $B(f, g)$ на декартовом произведении $F \times G$ локально выпуклых пространств $F$ и, допускающая представление вида



\[

B(f, g)=\sum_{i=1}^{\infty} \lambda_{i}\left\langle f, f_{i}^{\prime}\right\rangle\left\langle g, g_{i}^{\prime}\right\rangle,

\]



где $\left\{\lambda_{i}\right\}-$ суммируемая последовательность, $\left\{f_{i}^{\prime}\right\}$ и $\left\{g_{i}^{\prime}\right\}$ - равностепенно непрерывные последователь- ности в сопряженных к $F$ и $G$ пространствах $F^{\prime}$ и $G^{\prime}$ соответственно, а значение линейного функционала $a^{\prime}$ на векторе $a$ обозначается $\left\langle a, a^{\prime}\right\rangle$. Все Я. б. ф. непрерывны. Если $F$ - ядерное пространство, то для любого локально выпуклого пространства $G$ все непрерывные билинейные формы на $F \times G$ являются ядерными (т е ор ө м a o я д p e). Этот результат принадлежит А. Гротендику [1]; в приведенной форме теорема о ядре сформулирована в [2], другие формулировки см. в [3]. Справедливо и обратное утверждение: если для пространства $F$ выполняется заключение теоремы о ядре, то это пространство ядерно.



Для пространств гладких финитных функций теорему о ядре впервые получил Л. ІШварц [4]. Пусть $D$ - ядерное пространство всех бесконечно дифференцируемых функций с компактным носителем на прямой, наделенное стандартной локально выпуклой топологией ШІварца, так что сопряженное пространство $D^{\prime}$ состоит из всех обобщенных функций на прямой. Для частного случая $F=G=D$ теорема о ядре эквивалентна следующему утверждению: всякий непрерывный билинейный функционал на $D \times D$ имеет вид



$B(f, g)=\left\langle f\left(t_{1}\right) g\left(t_{2}\right), F\right\rangle=\int_{-\infty}^{\infty} F\left(t_{1}, t_{2}\right) f\left(t_{1}\right) g\left(t_{2}\right) d t_{1} d t_{2}$, где $f(t), g(t) \in D$ и $F=F\left(t_{1}, t_{2}\right)$ - обобщенная фуннкция от двух переменных. Аналогичную формулировку допускает теорема о ядре для пространств гладких финитных функций от нескольких переменных, пространств быстро убывающих функций и других конкретных ядерных пространств. Аналогичные результаты справедливы и для полилинейных форм.



Непрерывную билинеиную форму $B(f, g)$ на $D \times D$ можно отождествить с непрерывным линейным оператором $A: D \rightarrow D^{\prime}$ с помощыо равенства



\[

B(f, g)=\langle g, A f\rangle,

\]



тто приводит к формулировке т е о р е м ы II в а рц а о я д р е: для каждого непрерывного линейного отображения $A: D \rightarrow D^{\prime}$ существует такая однозначно определенная обобщенная функция $F\left(t_{1}, t_{2}\right)$ от двух переменных, что



\[

A: f\left(t_{1}\right) \mapsto \int_{-\infty}^{\infty} F\left(t_{1}, t_{2}\right) f\left(t_{2}\right) d t_{2}

\]



для всех $f \in D$. Другими словами, $A$ является интегральным оператором с ядром $F$.



Jum.: [1] G r othendieck A., Produits tensoriels topo$\log$ iques et espaces nucléaires, Providence, 1955; [2] П и ч А., Ядерные локально выпуклые пространства, пер. с нем., М., 1967 ; [3] Г ел ь ф а н д И. М., В и л е н к и н Н. Я., Некото-

рые применения гармонического анализа. Оснащенные гильбертовы пространства, M., 1961; [4] S c h w a r t z L., "Proc Intern. Congr. Math." (Camb., 1950), 1952 , v. 1, p. 220-30 [5] е г о ж e, «J. d'Analyse Math.», 1954/55, v. 4, p. 88-148. Г. ЈI. ЈІитвинов. $: A \rightarrow A^{\prime}$ алгебраич. систем- конгруәнция $\theta$ на алгебраич. системе $A$, состоящая из всех пар $(a, b) \in A \times A$, для к-рых $\varphi(a)=\varphi(b)$. Для всякой конгруәнции $\theta$ алгебраич. системы существует гомоморфизм $\varphi$ әтой системы, для к-рого $\theta$ является Я.к. Если $\theta-$ Я.к. сильного гомоморфиизма $\varphi$ алгебраич. системы $A$ на





\end{document}